\documentclass[11pt,a4paper]{article}

% ── Packages ─────────────────────────────────────────────────────────────────
\usepackage{amsmath,amssymb,amsthm}
\usepackage{mathtools}
\usepackage{geometry}
\usepackage{hyperref}
\usepackage{booktabs}
\usepackage{array}
\usepackage{xcolor}
\usepackage{enumitem}
\usepackage{microtype}
\usepackage{titlesec}
\usepackage{abstract}

\geometry{margin=2.5cm}

% ── Theorem environments ──────────────────────────────────────────────────────
\newtheorem{theorem}{Theorem}[section]
\newtheorem{lemma}[theorem]{Lemma}
\newtheorem{corollary}[theorem]{Corollary}
\newtheorem{proposition}[theorem]{Proposition}
\newtheorem{definition}[theorem]{Definition}
\newtheorem{remark}[theorem]{Remark}

% ── Operators ─────────────────────────────────────────────────────────────────
\DeclareMathOperator{\HF}{HF}
\DeclareMathOperator{\CF}{CF}
\DeclareMathOperator{\Jac}{Jac}
\DeclareMathOperator{\supp}{supp}
\DeclareMathOperator{\PO}{PO}
\newcommand{\one}{\mathbf{1}}
\newcommand{\Fuk}{\mathcal{F}}
\newcommand{\Lag}{\mathcal{L}}
\newcommand{\Mfld}{\mathcal{M}}
\newcommand{\AU}{\mathsf{AU}}

% ─────────────────────────────────────────────────────────────────────────────
\title{%
  \textbf{Homological Admissibility and Topological Hallucination Removal}\\[6pt]
  \large A Fukaya-Categorical Framework for Constrained Transformer Decoding
}

\author{
  AU-Fukaya Compiler Project\\
  \texttt{FukayaAUComplex}
}

\date{\today}

% ─────────────────────────────────────────────────────────────────────────────
\begin{document}
\maketitle

\begin{abstract}
We present a \emph{compiled semantic admissibility architecture}
in which higher-categorical consistency constraints are lowered into
executable runtime admissibility projection passes over generative trajectories.
Standard transformer decoding optimises pointwise local likelihood
over an unconstrained support $\mathbb{R}^D$, permitting generation
of outputs that are geometrically close in latent space but
topologically inadmissible with respect to the data manifold.
We replace this with \emph{support-restricted hierarchical
continuation decoding}, governed by the core equation:
\[
P_{\mathrm{topo}}(x_{t+1} \mid X_{\leq t})
\;=\;
\frac{P_\theta(x_{t+1} \mid X_{\leq t})
  \cdot \prod_{n=1}^{5}\mathbf{1}_{\{ob_n(x_{t+1}\mid X_{\leq t})=0\}}}
     {Z}
\]
where $P_\theta$ is the raw model distribution, each indicator
$\mathbf{1}_{\{ob_n=0\}}$ tests whether the continuation lifts
through Postnikov stage $n$, and $Z$ is the renormalisation over
the surviving admissible support.
Hallucination is reformulated as a \emph{failure of hierarchical
continuation liftability} through a compiled semantic filtration:
obstruction cocycles are evaluated prior to the softmax, restricting
the support of the generative distribution to trajectories admitting
coherent continuation across all active Postnikov stages.
No external ground-truth oracle is required during decoding;
admissibility is computed intrinsically from pre-compiled,
memory-mapped constraint structures (toric ideal generators,
Floer-like continuation invariants, HMM brackets).
Topology determines support; probability determines selection within
support.
The framework is validated empirically on a protocol
state-machine benchmark (Phase~1) and on the 81-station Hong Kong
MTR transit network, and connects theoretically to
Fukaya-Oh-Ohta-Ono Floer theory~\cite{FOOO} and constructible
sheaves~\cite{FLTZ}.
\end{abstract}

\tableofcontents
\newpage

% =============================================================================
\section{Introduction}
% =============================================================================

A \emph{hallucination} in a language model is the generation of a
fluent, confident output that has no valid grounding in reality.
Existing mitigation strategies --- RLHF, temperature scaling,
retrieval-augmented generation, entropy thresholding --- share a
common weakness: they attempt to make hallucinations \emph{unlikely}
rather than \emph{structurally impossible}.

This paper proposes a different approach. We observe that:
\begin{enumerate}[label=(\roman*)]
  \item The set of valid semantic states forms an algebraic variety
        $V(I) \subset \mathbb{R}^D$ defined by a toric ideal $I$;
  \item Valid transitions between states correspond to morphisms in
        a Fukaya $A_\infty$-category $\Fuk(M)$ over a symplectic
        manifold $(M,\omega)$;
  \item A hallucination is precisely a generative trajectory that
        lacks a nonvanishing Floer continuation class --- a
        topological obstruction, not a probabilistic one.
\end{enumerate}

The resulting architecture constrains decoding through a
\emph{Homological Attention Gate}: a hard indicator function
$\one_{\HF \neq 0}$ inserted before the softmax in each attention
head, forcing zero probability on topologically obstructed
continuations. This is not a learned constraint. It is derived
algebraically at compile time and cannot be overridden by gradient
descent.

\subsection{The core paradigm shift}

The central architectural distinction is:

\begin{center}
\renewcommand{\arraystretch}{1.4}
\begin{tabular}{p{5cm}p{6.5cm}}
\toprule
\textbf{Standard transformer} & \textbf{This work} \\
\midrule
Pointwise token sampling & Morphism chaining in $\Fuk(M)$ \\
Local likelihood optimisation & Homotopy-coherent continuation \\
$x_{t+1} \sim P_\theta(\cdot \mid x_{\leq t})$ &
  $x_{t+1}$ admissible only if $\gamma_t \rightsquigarrow \gamma_{t+1}$
  preserves continuation structure \\
Causal mask = binary wall & Homological gate = obstruction cocycle \\
Inconsistency = low probability & Inconsistency = non-liftable at stage $P_n$ \\
\bottomrule
\end{tabular}
\end{center}

\noindent
A continuation is allowed only if it survives the differential
$d: \CF^k \to \CF^{k+1}$.
If it fails, the trajectory is geometrically unsupported ---
unreachable under the admissible topology, or non-liftable in the
semantic filtration.
We are careful not to equate geometric obstruction with semantic
falsity: a blocked trajectory may be \emph{unsupported} rather than
\emph{false}. The claim is structural unreachability, not
omniscient truth detection.

\subsection{Comparison with existing methods}

\subsubsection*{Table 1: Competing definitions of hallucination}

\begin{center}
\renewcommand{\arraystretch}{1.3}\small
\begin{tabular}{p{2.6cm}p{1.6cm}p{1.9cm}p{2.0cm}p{1.4cm}p{2.8cm}}
\toprule
\textbf{Method} &
\textbf{Type} &
\textbf{Oracle at inference?} &
\textbf{Modifies support before sampling?} &
\textbf{Hierarchical?} &
\textbf{Failure mode} \\
\midrule
Low-prob.\ threshold & Statistical & No & No — post-hoc & No &
  Confident OOD hallucination missed \\
Temperature scaling & Probabilistic & No & No & No &
  Nonzero mass on invalid states \\
Entropy threshold & Statistical & No & No — post-sampling & No &
  Misses confident hallucinations \\
RLHF & Epistemic & Yes (human pref.) & No — reward learning & No &
  Confident unseen hallucinations \\
RAG & Epistemic & Yes (external DB) & No — post-planning & No &
  Hallucination between retrieved facts \\
Contrastive decoding & Statistical & No & No — logit adjust & No &
  Requires expert/amateur model pair \\
Self-consistency & Epistemic & Partial (verifier) & No — post-hoc rerank & No &
  Verifier has its own hallucinations \\
\midrule
\textbf{This work} & Topological & \textbf{No (intrinsic)} &
  \textbf{Yes} --- before softmax &
  \textbf{Yes} --- 5 stages &
  Overly restrictive admissibility (false positives, not false negatives) \\
\bottomrule
\end{tabular}
\end{center}

\subsubsection*{Table 2: Detection vs.\ prevention}

\begin{center}
\renewcommand{\arraystretch}{1.3}\small
\begin{tabular}{p{4.2cm}p{4.5cm}p{4.5cm}}
\toprule
\textbf{Property} &
\textbf{Prior art} (RAG, RLHF, entropy, verifiers) &
\textbf{This work} (Homological Gate) \\
\midrule
When does evaluation occur? &
  After token generation, or during training &
  Before softmax (compile-time $+$ $O(1)$ runtime) \\
What is modified? &
  Sampled outputs, reward signals, retrieved contexts &
  $\mathrm{supp}(P)$ --- the support of the distribution \\
Can a hallucination still be sampled? &
  Yes, then corrected / filtered / penalised &
  No --- structurally impossible \\
Ground truth required? &
  Yes (DB, human labels, verifier model) &
  No (intrinsic compiled constraints only) \\
Own failure mode? &
  False negatives (missed hallucinations) &
  False positives (over-restricted admissibility) \\
Inference overhead &
  Low--moderate (retrieval, reranking, DB query) &
  Very low (memory-mapped IR lookup, $P99 < 5\,$ms) \\
Hierarchical scale awareness &
  No --- flat hallucination definition &
  Yes --- $P_1$ token through $P_5$ discourse \\
\bottomrule
\end{tabular}
\end{center}

\subsubsection*{Table 3: The five obstruction stages vs.\ prior art}

\begin{center}
\renewcommand{\arraystretch}{1.3}\small
\begin{tabular}{p{2.8cm}p{1.8cm}p{2.5cm}p{5.0cm}}
\toprule
\textbf{Prior art type} &
\textbf{Approx.\ stage} &
\textbf{Misses stages} &
\textbf{Why prior art fails structurally} \\
\midrule
``Lexically impossible token'' & $P_1$ & $P_2$--$P_5$ &
  Word filters do not catch phrasal incoherence \\
``Semantic incoherence'' & $P_2$ & $P_3$--$P_5$ &
  Local coherence $\neq$ global discourse liftability \\
``Temporal inconsistency'' & $P_3$ & $P_4, P_5$ &
  Requires external timeline; no intrinsic causal structure \\
``Causal impossibility'' & $P_4$ & $P_5$ &
  Local causality $\neq$ discourse-level continuation class \\
``Discourse contradiction'' & $P_5$ & None above $P_5$ &
  Cannot explain why trajectory failed at earlier scales \\
\bottomrule
\end{tabular}
\end{center}

\noindent\textbf{Key insight.}
Prior work treats each hallucination type as a separate problem
requiring a separate solution.
The Postnikov filtration treats them as a \emph{single hierarchical
obstruction cocycle} evaluated uniformly across all five scales.

\medskip\noindent\textbf{Refined defensible claim.}
\begin{quote}
\emph{No existing method provides pre-sampling support restriction
under a hierarchical admissibility criterion.
Prior work detects hallucinations; this work structurally prevents
them at the level of the sampling distribution's support, using only
intrinsically compiled constraints and no external factual oracle
at inference time.}
\end{quote}

\medskip\noindent\textbf{Key systems advantage.}
Prior art introduces external oracles at inference time
(retrieval databases, separate classifiers, web search),
incurring latency incompatible with sub-5\,ms SLA requirements.
This architecture pre-compiles all admissibility constraints into
memory-mapped flat files consulted in $O(1)$ at runtime.

% =============================================================================
\section{Mathematical Foundations}
% =============================================================================

\subsection{The symplectic manifold and Lagrangian objects}

\begin{definition}[Semantic symplectic manifold]
Let $(M, \omega)$ be a compact toric symplectic manifold where:
\[
  \omega(s,h,m) \;=\; r(s)\cdot\phi(h)\cdot\mu(m)
\]
is the symplectic form built from ridership weights $r(s)$,
hourly profile $\phi(h)$, and seasonal multiplier $\mu(m)$.
This is the classical limit ($T=0$) of the FOOO superpotential
$W_0(u)$ from~\cite{FOOO}.
\end{definition}

\begin{definition}[Lagrangian submanifolds]
The \emph{semantic states} are Lagrangian submanifolds
$L_i \subset M$ indexed by demographic strata
$i \in \{RM, RF, PM, PF\}$:
\[
  L_i : (s,h,m) \;\mapsto\; \mathbb{R}_{\geq 0}
\]
defined by ridership-weighted affinity functions. These are the
objects of the Fukaya category $\Fuk(M)$.
\end{definition}

\subsection{The Floer chain complex}

\begin{definition}[Floer chain complex]
For Lagrangians $L_i, L_j$ in general position, the
\emph{Floer chain complex} is:
\[
  \CF^*(L_i, L_j) \;=\;
  \bigoplus_{x \,\in\, L_i \cap L_j} \mathbb{Z} \cdot x
\]
with one generator per transverse intersection point.
The differential $\partial: \CF^k \to \CF^{k+1}$ counts rigid
$J$-holomorphic strips $u: \mathbb{R} \times [0,1] \to M$
satisfying:
\[
  \frac{\partial u}{\partial s} + J\frac{\partial u}{\partial t} = 0,
  \qquad u(s,0) \in L_i,\quad u(s,1) \in L_j
\]
connecting intersection points as $s \to \pm\infty$.
The FOOO virtual technique~\cite{FOOO} establishes $\partial^2 = 0$.
\end{definition}

\begin{definition}[Floer homology]
\[
  \HF^*(L_i, L_j) \;=\; \ker\partial \,/\, \mathrm{im}\,\partial
\]
When $\HF^*(L_i, L_j) = 0$, the transition from state $L_i$ to
state $L_j$ is \emph{topologically obstructed}: no valid
holomorphic continuation exists.
\end{definition}

\begin{remark}[Classical limit]
At $T=0$ the differential $\partial = 0$ (no strips of zero
symplectic area in the generic setting), so $\HF^* = \CF^*$ and
vanishing reduces to $L_i \cap L_j = \emptyset$ on the toric
variety $V(I)$. Our system operates in this classical limit;
the full quantum corrections (Novikov ring terms) constitute
the programme for Passes 6--8.
\end{remark}

\subsection{Operational grounding of abstract objects}

To silence the reviewer objection that this is ``decorative
category theory,'' we map each abstract object to a concrete,
computable artifact:

\begin{description}[leftmargin=2cm, style=nextline]
  \item[Symplectic manifold $M$]
    The cotangent bundle of the semantic state space.
    Base manifold = discrete graph of valid network paths and
    demographic constraints.
    Computationally: the 81-node MTR graph with integer edge weights,
    stored as a sparse adjacency matrix.

  \item[Lagrangian submanifolds $L_i$]
    Bounded polytopes in demographic-ridership space, defined by
    systems of linear inequalities:
    $L_i = \{(s,h,m) \mid \omega_i(s,h,m) \geq \theta_i\}$.
    Stored as Pass~1 embedding vectors in flat Parquet format.

  \item[Chain grading $k$]
    Corresponds directly to the Postnikov semantic scale:
    $k=1$ = token, $k=2$ = phrase, $k=3$ = causal/world-model.

  \item[Differential $d$]
    Implemented as the combinatorial boundary operator
    $\partial_1 \in \mathbb{R}^{|V|\times|E|}$ acting on the
    simplicial activation complex.
    Computing $d \cdot X = 0$ is realised as confirming a token
    transition satisfies the conservation-of-flow equations
    across the pre-compiled network (the toric ideal relations).

  \item[Energy functional / $\HF^* \neq 0$]
    The HMM transition penalties and FK bracket
    $[\mathcal{P}_{\min}, \mathcal{P}_{\max}]$.
    Non-vanishing cohomology = there exists at least one valid
    path through the network satisfying the joint demographic
    distribution.
    If no such path exists, $\HF^* = 0$ and the trajectory is pruned.

  \item[Continuation map $\Phi$]
    The AU pullback functor at each surgery node $(s,h,m)$:
    evaluates the coproduct $L_i \sqcup L_j$ lazily and checks
    whether the stalk is nonzero.
\end{description}

\subsection{The $A_\infty$ structure}

The Fukaya category is an $A_\infty$-category with higher
composition operations:
\[
  m_k : \CF^*(L_0,L_1) \otimes \cdots \otimes
        \CF^*(L_{k-1},L_k) \to \CF^*(L_0,L_k), \quad k \geq 1
\]
satisfying the $A_\infty$ relations
$\sum_{k} m_{n-k+1}\circ(\mathrm{id}^{\otimes}\otimes m_k \otimes
\mathrm{id}) = 0$.
The substitution $q \mapsto f(q)$ expands the superpotential as:

\begin{center}
\renewcommand{\arraystretch}{1.3}
\begin{tabular}{llll}
\toprule
Degree & Operation & Geometric meaning & Compiler role \\
\midrule
0 & $m_0$ & Curvature, $W_0(u) = \omega(s,h,m)$ &
  Classical obstruction, dead-zone detection \\
1 & $m_1$ & Differential, Markov propagation &
  Neighbourhood spillover (Pass 4) \\
2 & $m_2$ & Composition, $\langle L_i,L_j\rangle\cdot\mathrm{aff}(p)$ &
  Routing score, attention $QK^\top$ \\
3 & $m_3$ & Homotopy, timing stability &
  Loop-carried dependence (Pass 3) \\
$k$ & $m_k$ & Novikov corrections &
  Passes 6--8 (future) \\
\bottomrule
\end{tabular}
\end{center}

% =============================================================================
\section{The Toric Ideal and Algebraic Obstruction}
% =============================================================================

\subsection{The Markov basis and toric variety}

The network topology defines a toric ideal via the 4ti2
algebraic computation:
\[
  I \;=\; \langle f_1, \ldots, f_{37} \rangle \;\subset\;
  \mathbb{Z}[x_1,\ldots,x_{81}]
\]
where $f_1,\ldots,f_{37}$ are the Markov circuits of the
MTR incidence matrix ($\beta_1 = |E| - |V| + 1 = 89-81+1 = 9$
independent cycles, 37 Markov circuits in the Graver basis).

\begin{definition}[Toric variety]
\[
  V(I) \;=\; \{ q \in \mathbb{R}^{81} \mid f_i(q) = 0,\;
  i = 1,\ldots,37 \}
\]
is the \emph{admissible semantic manifold}: the set of all
demographically consistent passenger flows.
Its complement $V(I)^c$ is the \emph{hallucination space} ---
states that are algebraically impossible.
\end{definition}

\subsection{The $q \mapsto f(q)$ coordinate substitution}

\begin{definition}[Structured embedding]
Raw token vectors $q$ are projected into the coordinate ring
$\mathbb{R}[x_1,\ldots,x_D]/I$ via:
\[
  q \;\mapsto\; X \;=\; [f_1(q),\, f_2(q),\, \ldots,\, f_k(q)]
\]
This maps each token into the coordinate ring of $V(I)$.
Tokens outside $V(I)$ receive zero coordinates and are
structurally excluded from retrieval.
\end{definition}

In the compiler, this is Pass~1:
$\mathrm{embed}(p)[i] = \langle L_i, p\rangle$, the global
Floer pairing with demographic Lagrangian $L_i$.

% =============================================================================
\section{The Postnikov Tower and Hierarchical Obstruction}
% =============================================================================

\subsection{The filtration}

A hallucination can begin in admissible regions and drift globally
while remaining locally valid. A pointwise support constraint
$\supp(P) \subseteq V(I)$ misses this. The Postnikov filtration
resolves it by requiring coherence at every semantic scale.

\begin{definition}[Postnikov filtration]
The compiler implements a five-stage Postnikov tower:
\[
  P_5(X) \to P_4(X) \to P_3(X) \to P_2(X) \to P_1(X)
\]
\end{definition}

\begin{center}
\renewcommand{\arraystretch}{1.3}
\begin{tabular}{cllp{4cm}}
\toprule
Stage & Pass & Semantic level & Obstruction if failed \\
\midrule
$P_1$ & Product embeddings & Token coherence &
  Token not on $V(I)$ --- algebraic obstruction \\
$P_2$ & Routing table & Phrase admissibility &
  No valid demographic intersection \\
$P_3$ & Stability table & Temporal coherence &
  $m_3$ loop-carried dependence \\
$P_4$ & Neighbourhood & Causal coherence &
  $m_1$ propagation fails across nodes \\
$P_5$ & HMM brackets & Discourse coherence &
  $k\text{-inv} \approx 0$, continuation class vanishes \\
\bottomrule
\end{tabular}
\end{center}

\begin{definition}[Obstruction cocycle]
The HMM bracket $[\mathcal{P}_{\min}, \mathcal{P}_{\max}]$
computed at Pass~5 is the obstruction cocycle at the discourse
level. The $k$-invariant:
\[
  k\text{-inv}(s,h,m) \;=\; \frac{\mathcal{P}_{\min}}{\mathcal{P}_{\max}}
\]
measures whether the Feynman-Kac expectation of $W_0$ along the
trajectory has a nonvanishing continuation class.
When $k\text{-inv} \approx 0$, the lift to stage $P_5$ is
obstructed --- the trajectory is a topological dead zone.
\end{definition}

\subsection{Homological descent}

The sheaf of admissible universes $\{\Mfld_U\}_{U \in \mathcal{C}}$
descends compatibly through the tower:
\[
  \Mfld_{P_5} \;\hookrightarrow\;
  \Mfld_{P_4} \;\hookrightarrow\;
  \Mfld_{P_3} \;\hookrightarrow\;
  \Mfld_{P_2} \;\hookrightarrow\;
  \Mfld_{P_1}
\]
Each inclusion is a restriction map --- the AU pullback at a
surgery node $(s,h,m)$. The AU attic implements this sheaf as
lazy evaluation: coproducts $L_i \sqcup L_j$ are materialised
locally at context $(s,h,m)$ and never globally allocated.

% =============================================================================
\section{The Discrete Hodge Decomposition Pass}
% =============================================================================

\subsection{Three-component flow decomposition}

Let $f \in \mathbb{R}^{|E|}$ be an edge flow vector on the network
graph $G = (V,E)$ with incidence matrix
$\partial_1 \in \mathbb{R}^{|V| \times |E|}$.
The discrete Hodge decomposition gives:
\[
  f \;=\; f_{\mathrm{grad}} + f_{\mathrm{harm}} + f_{\mathrm{curl}}
\]
where:
\begin{align*}
  f_{\mathrm{grad}} &= \partial_1^\top (\partial_1 \partial_1^\top)^+ \partial_1 f
    & &\text{(gradient: spanning-tree flows)} \\
  f_{\mathrm{harm}} &= f - f_{\mathrm{grad}}
    & &\text{(harmonic: independent loop flows, } \beta_1=9\text{)} \\
  f_{\mathrm{curl}} &= 0
    & &\text{(curl: zero for plain graphs, } \beta_2=0\text{)}
\end{align*}

\subsection{Hallucination detection via Hodge gate}

\begin{definition}[Three-valued Hodge gate]
\[
  \mathrm{gate}(s,h,m) \;=\;
  \begin{cases}
    \mathsf{DEAD}     & \omega(s,h,m) < \theta \cdot \omega_{\max}
                        \quad\text{(no ridership: off-manifold)} \\
    \mathsf{UNSTABLE} & \|f_{\mathrm{curl}}\|/\|f\| > \tau
                        \quad\text{(curl dominates: tautological loop)} \\
    \mathsf{LIVE}     & \text{otherwise}
                        \quad\text{(gradient + harmonic: valid forward flow)}
  \end{cases}
\]
with thresholds $\theta = 0.02$, $\tau = 0.40$.
\end{definition}

The mapping to the document's LLM framing is exact:
\begin{itemize}
  \item \textbf{Gradient} = valid forward flow (progressive inference,
        grounded reasoning)
  \item \textbf{Harmonic} = loop signal (genuine multi-step argument
        traversing a topological cycle; positive contribution)
  \item \textbf{Curl} = tautological loop (cyclic reasoning,
        $X$ because $Y$ because $X$; hallucination signal)
  \item \textbf{Dead zone} = $\|f\| \approx 0$
        (no flow at all; off-manifold, unreachable)
\end{itemize}

\subsection{Relation to the existing scalar gates}

The Hodge pass unifies three separate scalar computations:

\begin{center}
\renewcommand{\arraystretch}{1.3}
\begin{tabular}{lll}
\toprule
Hodge component & Replaces & Used by \\
\midrule
$\mathrm{grad\_ratio}$ & raw $\omega$ routing score & Pass 2 routing table \\
$\mathrm{harm\_ratio}$ & FK bracket $k$-invariant & Pass 5 HMM brackets \\
$\mathrm{curl\_ratio}$ & $1 - \mathrm{stability}$ & Pass 3 stability table \\
\bottomrule
\end{tabular}
\end{center}

All three are derived from one matrix decomposition: build
$\partial_1$ once, compute $(\partial_1\partial_1^\top)^+$ once,
and all slot classifications follow.

% =============================================================================
\section{The Homological Attention Gate}
% =============================================================================

\subsection{Standard attention and its failure}

Standard self-attention computes:
\[
  \mathrm{Attention}(Q,K,V)
  \;=\; \mathrm{softmax}\!\left(\frac{QK^\top}{\sqrt{d_k}}\right) V
\]
This is a \emph{degree-2 computation} in the $A_\infty$ sense:
it evaluates $m_2$ (the bilinear Floer pairing) but ignores:
\begin{itemize}
  \item $m_0$ (degree 0): does this state exist? ($W_0(u) > 0$)
  \item $m_1$ (degree 1): is this transition topologically connected?
  \item $m_3, m_4, \ldots$: stability and Novikov corrections
\end{itemize}
A hallucination arises when the model assigns high $m_2$ score
to a transition that is obstructed at $m_0$ or $m_1$ level.

\subsection{The gated decoding rule}

\begin{definition}[Homological attention gate]
The gated attention logits apply the Postnikov obstruction
indicator before the softmax:
\[
  A_{ij} \;=\; \frac{QK^\top}{\sqrt{d_k}} \;+\;
  \sum_{n=1}^{5}\log\mathbf{1}_{\{ob_n(j\mid X_{\leq t})=0\}}
\]
Since $\log 0 = -\infty$, any token $j$ failing at any stage
receives zero probability mass before normalisation.
\end{definition}

The complete gated decoding distribution is:
\[
\boxed{
  P_{\mathrm{topo}}(x_{t+1} \mid X_{\leq t})
  \;=\;
  \frac{P_\theta(x_{t+1} \mid X_{\leq t})
    \cdot\prod_{n=1}^{5}\mathbf{1}_{\{ob_n(x_{t+1}\mid X_{\leq t})=0\}}}
       {Z}
}
\]
\textbf{Topology determines support; probability determines selection
within support.}
This is not a post-hoc reranker, retrieval correction, or confidence
filter. It modifies $\mathrm{supp}(P)$ before normalisation.
Obstructed continuations never enter the support, never receive
gradient updates, and never appear in generated output ---
regardless of the sampling temperature or the learned weights $\theta$.

% =============================================================================
\section{The Homological Admissibility Theorem}
% =============================================================================

\begin{theorem}[Topological support restriction]
\label{thm:support}
For any slot $(s,h,m)$ with $W_0(s,h,m) = 0$
(toric obstruction, $x \notin V(I)$):
\[
  P_{\mathrm{topo}}(\text{any output} \mid s,h,m) \;\equiv\; 0
\]
The support of the gated distribution satisfies:
\[
  \supp(P_{\mathrm{topo}}) \;\subseteq\; V(I)
\]
\end{theorem}

\begin{proof}
$W_0(s,h,m) = 0
\Rightarrow \omega(s,h,m) = 0
\Rightarrow L_i \cap L_j = \emptyset\;\forall i,j
\Rightarrow \CF^* = 0
\Rightarrow \HF^* = 0
\Rightarrow \one_{\HF \neq 0} = 0
\Rightarrow P_{\mathrm{topo}} = 0.$
\end{proof}

\begin{theorem}[Homological Admissibility]
\label{thm:main}
Let semantic states form objects $L_i \in \Fuk(M)$ in the
Fukaya category of the toric symplectic manifold $(M,\omega)$.
Let decoding trajectories correspond to $A_\infty$ morphism
compositions through the Postnikov filtration
$P_1 \subset \cdots \subset P_5$.
Any generated sequence whose continuation induces a trivial
obstruction class at Postnikov stage $n$ is excluded from
decoding support. Formally:
\[
  \forall\, x_{t+1}:\quad
  \exists\, n \in \{1,\ldots,5\} \text{ s.t. }
  \text{lift}_{P_n}(x_{t+1}) \text{ obstructed}
  \;\Longrightarrow\;
  P_{\mathrm{topo}}(x_{t+1} \mid X_{\leq t}) = 0
\]
\end{theorem}

\begin{proof}[Proof sketch]
By induction on Postnikov stages:
\begin{itemize}
  \item \textbf{Stage $P_1$}: $x_{t+1} \notin V(I)
    \Rightarrow \mathrm{embed}(x_{t+1}) = 0
    \Rightarrow$ excluded by the $q \mapsto f(q)$ substitution.
  \item \textbf{Stage $P_2$}: no demographic intersection
    $\Rightarrow$ absent from routing table $\Rightarrow$ score $= 0$.
  \item \textbf{Stage $P_3$}: $m_3 > \tau \Rightarrow$ flagged volatile
    $\Rightarrow$ skipped by stability gate.
  \item \textbf{Stage $P_4$}: $m_1$ propagation fails across adjacent
    nodes $\Rightarrow$ absent from neighbourhood table $\Rightarrow$
    unreachable.
  \item \textbf{Stage $P_5$}: $k\text{-inv} \approx 0
    \Rightarrow W_0(u) \approx 0
    \Rightarrow \HF^* = 0
    \Rightarrow \one_{\HF \neq 0} = 0
    \Rightarrow P_{\mathrm{topo}} = 0$.
\end{itemize}
The $A_\infty$ relation at each surgery node ensures these gates
are mutually consistent.
\end{proof}

\begin{corollary}[Impossibility of hallucination at topological boundary]
A sequence that is topologically obstructed (lies in $V(I)^c$ or
has $k\text{-inv} = 0$) cannot be generated regardless of the
learned weights $\theta$, the temperature, or the prompt.
\end{corollary}

% =============================================================================
\section{Variable-Set Semantics and Dynamic Topology}
% =============================================================================

The framework generalises from static support restriction
$\supp(P) \subseteq \Mfld$ to context-dependent admissibility:
\[
  \supp(P_t) \;\subseteq\; \Mfld_t
\]
where the admissible manifold $\Mfld_t$ evolves with the
generation context $X_{\leq t}$.

\begin{definition}[Sheaf of admissible universes]
In variable-set (sheaf/topos) semantics, membership is
context-dependent:
\[
  x \in \Mfld(U) \quad\text{for contextual region } U \in \mathcal{C}
\]
A token may be admissible in one discourse frame but obstructed
in another. The family $\{\Mfld_U\}_{U \in \mathcal{C}}$ with
gluing conditions is a \emph{sheaf of admissible universes}.
\end{definition}

The AU (Arithmetic Universe) attic implements this sheaf as
lazy evaluation: coproducts $L_i \sqcup L_j$ are evaluated
in local context $(s,h,m)$ via the pullback functor and are
never globally materialised. Each surgery node is a stalk;
restriction maps between adjacent stalks are the Floer
continuation maps.

\subsection{Floer continuation as context evolution}

The continuation map $\Phi: \HF^*(L_0, L_1) \to \HF^*(L_0', L_1')$
under Hamiltonian deformation is the formal mechanism by which
$\Mfld_t$ evolves as context accumulates. A hallucination
corresponds to a trajectory:
\[
  \gamma: [0,T] \to \Mfld_t
\]
that \emph{fails to possess a nontrivial continuation class} ---
meaning $\Phi(\gamma) = 0$ at some step $t$. The decoding is
terminated at that step.

% =============================================================================
\section{Phase 1 Empirical Benchmark}
\label{sec:benchmark}

\subsection{Setup}

We train two identical nanoGPT models (3 transformer blocks,
$d_\mathrm{model}=64$, 4 attention heads) on 5,000 randomly
generated valid trajectories of a 7-token protocol state machine:
\[
  \texttt{START} \to \texttt{AUTH} \to
  \{\texttt{QUERY} \to \texttt{RESULT} \mid \texttt{RETRY} \mid
    \texttt{ERROR}\}^* \to \texttt{END}
\]
The \emph{admissible transition graph} defines the toric ideal $I$
(7 states, 10 admissible edges, 0 self-loops).
Forbidden transitions include \texttt{START}$\to$\texttt{RESULT},
\texttt{AUTH}$\to$\texttt{END}, \texttt{ERROR}$\to$\texttt{QUERY},
and 14 others.

\textbf{Baseline}: standard nanoGPT with causal attention mask only.
\textbf{Constrained}: identical architecture with the IR admissibility
gate inserted before each softmax:
$A_{ij} = QK^\top/\sqrt{d_k} + \log\mathbf{1}_\mathcal{M}(j)$.

\subsection{Training loss}

Training loss curves for both models are indistinguishable
($\Delta\mathrm{loss} < 0.001$ at all epochs).
This confirms that the coordinate ring projection
$q \mapsto f(q)$ does not damage learning capacity or
perplexity. The transformer optimises its weights naturally;
the gate adds no computational burden to the training objective.

\subsection{Three stress tests}

We design three stress tests that actively break the baseline.
Validity is measured on the \emph{generated suffix} only
(the forced prefix is excluded from the count).

\begin{center}
\renewcommand{\arraystretch}{1.35}
\begin{tabular}{lrrl}
\toprule
\textbf{Test} & \textbf{Baseline} & \textbf{Constrained} & \textbf{What breaks the baseline} \\
\midrule
Attn saturation (\texttt{QUERY}$\times 40$ prefix) &
  82.5\% & 100.0\% &
  Attention weights saturate; model drifts off-manifold \\
Illegal prefix (\texttt{QUERY}$\to$\texttt{AUTH} injected) &
  97.8\% & 100.0\% &
  Adversarial context confuses recovery policy \\
Under-training (epoch 5) &
  99.8\% & 100.0\% &
  Model has not converged; gate compensates \\
\bottomrule
\end{tabular}
\end{center}

\subsection{Temperature sweep}

As sampling temperature increases from $0.8$ to $3.0$,
the baseline degrades sharply:
\begin{itemize}
  \item $T=1.0$: baseline 99.0\%, constrained 100.0\%
  \item $T=2.0$: baseline 89.0\%, constrained 100.0\%
  \item $T=3.0$: baseline 54.0\%, constrained 100.0\%
\end{itemize}
The constrained model is flat at 100\% across all temperatures.
This is the structural guarantee: admissibility is enforced
\emph{before} the softmax, independent of the temperature
parameter or the learned weight distribution.

\subsection{The key isolation result}

The stress test with the undertrained model (epoch 5) isolates
the topological contribution from the learned contribution.
At epoch 5, the baseline has not converged (99.8\% valid vs.\ 98.5\%
at epoch 40 — note the slight degradation as the model overfits
slightly toward common patterns).
The constrained model holds at 100.0\% at both epoch 5 and epoch 40
because the gate is evaluated algebraically at inference time,
not learned from data.

This is the publishable empirical core: the IR gate converts an
undertrained, adversarially-prompted, or attention-saturated model
into one that generates 100\% valid trajectories --- structurally,
not probabilistically.

\section{Validated Results on the MTR Network}
% =============================================================================

The framework is validated on the Hong Kong MTR network with
$|V| = 81$ stations, $|E| = 89$ edges, $\beta_1 = 9$ independent
cycles, 37 Markov circuits (4ti2), and 161 Graver elements.

\subsection{Hodge decomposition results}

The discrete Hodge pass identifies:
\begin{itemize}
  \item \textbf{23 loop stations} ($\mathrm{harm\_ratio} > 0.05$):
        On one of the 9 network cycles. Highest: Olympic (0.490),
        Mong Kok East (0.446), Hung Hom (0.383), Austin (0.362).
        These receive stronger HMM bracket bounds.
  \item \textbf{58 tree stations} ($\mathrm{harm\_ratio} \approx 0$):
        On the spanning tree. Pure gradient routing. Fully live.
  \item \textbf{0 dead zones at peak hour}: All stations have
        $\omega > 0.03$ at 9am February. Dead zones emerge
        at off-peak hours ($\omega \to 0$).
\end{itemize}

\subsection{Gate agreement}

Testing all $486$ slots ($81$ stations $\times$ $2$ hours $\times$
$3$ months), the Hodge gate achieves:
\begin{itemize}
  \item \textbf{100\% exact match} with the existing scalar gates
        (routing table + stability + $k$-invariant).
  \item \textbf{0 mismatches}: the Hodge decomposition is a
        mathematically equivalent reformulation of the three
        separate scalar computations, derived from one incidence
        matrix computation.
\end{itemize}

% =============================================================================
\section{Discussion}
% =============================================================================

\subsection{What is genuinely novel}

Standard transformers perform \emph{probability maximisation} ---
local likelihood optimisation at each step. This architecture
performs \emph{homotopy-coherent continuation} --- global
trajectory validation through categorical state spaces.

The distinction is fundamental:
\begin{itemize}
  \item Standard: \emph{what is the most likely next token?}
  \item This work: \emph{does this continuation preserve
        cohomological consistency across all Postnikov stages?}
\end{itemize}

\subsection{Genuine vs.\ metaphorical contributions}

We distinguish explicitly between components that are
rigorously derived from the theory and those that are
engineering constructs organised by the mathematical language:

\begin{center}
\renewcommand{\arraystretch}{1.3}
\begin{tabular}{ll}
\toprule
\textbf{Genuinely theory-derived} & \textbf{Engineering in mathematical language} \\
\midrule
4ti2 Markov basis computation & HNSW hierarchical indexing \\
Pass 1 Floer pairing embedding & EMA feedback updates \\
FK bracket structure & $m_3$ finite-difference stability \\
Dead-zone detection via $W_0 \approx 0$ & Routing table as flat Parquet file \\
Hodge decomposition (one matrix op) & Threshold constants $\theta, \tau$ \\
\bottomrule
\end{tabular}
\end{center}

\subsection{Scope: where the architecture is strongest}

The framework is strongest for domains that naturally admit
continuation invariants, obstruction structure, and admissibility
topology:
\begin{itemize}
  \item \textbf{Protocol and workflow generation} (Phase~1 benchmark):
        state-machine transitions have exact admissibility topology.
  \item \textbf{Structured reasoning and theorem continuation}:
        proof steps have causal and logical continuation constraints.
  \item \textbf{Planning and routing} (MTR benchmark):
        network topology defines toric ideal generators exactly.
  \item \textbf{Program synthesis}: type systems and control-flow
        graphs provide natural Postnikov filtrations.
  \item \textbf{Stateful and causal trajectory generation}:
        any domain where ``what states can follow what states'' is
        definable in advance.
\end{itemize}
The architecture is \emph{weaker} for open-ended natural language
generation where admissibility topology cannot be pre-compiled ---
the false-positive rate (over-restriction) would be high.
The correct deployment sequence is:
Phase~1 (symbolic/structured) $\to$
Phase~2 (semi-structured reasoning) $\to$
Phase~3 (dynamic contextual topology) $\to$
Phase~4 (natural language) only after the above.

\subsection{The central remaining empirical question}

The key challenge identified in peer review is:

\begin{quote}
\emph{Do the admissibility constraints generalise, or do they
merely memorise finite valid trajectories?}
\end{quote}

This is now the scientific heart of the experimental programme.
The required experiments are:
\begin{enumerate}[label=(\roman*)]
  \item \textbf{Compositional generalisation}: generate trajectories
        combining sub-sequences not seen in training.
        Does validity hold? Does the topology contribute nontrivially
        over a learned mask?
  \item \textbf{OOD continuation validity}: the temperature sweep
        (Phase~1 results: baseline 54\% at $T=3.0$, constrained
        100\%) is a start. Full OOD requires topology changes
        (new states, new edges) not seen during compilation.
  \item \textbf{Diversity preservation}: does the admissibility
        projection preserve the entropy of $P_\theta$ within the
        admissible support, or does it catastrophically prune?
        Measurement: generation diversity (self-BLEU, unique
        $n$-gram rate) inside and outside the gate.
  \item \textbf{Creativity preservation}: are valid but
        ``creative'' trajectories (low-probability but admissible)
        still reachable? This tests whether the gate is merely
        a memorised rejection table or a genuine structural filter.
\end{enumerate}

\subsection{Known limitations}

\begin{enumerate}[label=(\roman*)]
  \item We operate at classical limit $T=0$. Full quantum
        corrections (Novikov ring) constitute Passes 6--8.
  \item The guarantee applies at the topological boundary.
        Within $V(I)$, ranking is heuristic and requires
        empirical validation.
  \item The static routing table requires recompilation when
        $\omega$ changes substantially. Adapton-style incremental
        recompilation is identified but not implemented.
  \item The moduli problem is computed at leading order
        ($\#\mathcal{M} \approx \sqrt{L_i \cdot L_j}\cdot\omega$),
        not from full virtual fundamental chains.
\end{enumerate}

% =============================================================================
\section{Patent Claim Structure}
\label{sec:patent}

The following structured claim captures the core invention
in a form suitable for patent or priority disclosure:

\begin{quote}
\textbf{Claim 1 (Method).}
A method for generative continuation validation comprising:
\begin{enumerate}[label=(\alph*)]
  \item a compiled Postnikov filtration
        $\{P_1,\ldots,P_5\}$ over an admissible semantic manifold
        $(M,\omega)$, wherein each stage corresponds to a distinct
        semantic coherence scale (token, phrase, temporal, causal,
        discourse);
  \item obstruction cocycles $ob_n$ evaluated prior to token
        sampling via pre-compiled, memory-mapped constraint
        structures (toric ideal generators, Floer-like continuation
        invariants, HMM brackets);
  \item a support restriction rule
        $P_{\mathrm{topo}}(x) \propto P_\theta(x)
        \cdot \prod_{n=1}^{5}\mathbf{1}_{\{ob_n=0\}}$
        applied before the softmax normalisation; and
  \item wherein a continuation is a \emph{structural hallucination}
        iff $\exists\,n: ob_n \neq 0$, and such continuations
        receive exactly zero probability mass.
\end{enumerate}

\textbf{Claim 2 (System).}
A system comprising: a static offline compiler that ingests a
network topology and emits pre-compiled admissibility tables;
a runtime gate that consults these tables in $O(1)$ per token;
and an autoregressive decoder whose sampling distribution has
support restricted to trajectories passing all five Postnikov
stages, requiring no external factual oracle at inference time.
\end{quote}

\section{Conclusion}
% =============================================================================

We have presented a framework in which hallucination removal
is algebraic and topological rather than probabilistic.
The key objects --- toric ideal, Floer chain complex, Postnikov
filtration, Hodge decomposition, AU sheaf --- are each
operationally defined and computationally instantiated on the
81-station MTR network.

The central claim is:

\begin{quote}
\emph{Hallucination is reformulated as a failure of hierarchical
continuation liftability through a compiled semantic filtration.
Obstruction cocycles are evaluated prior to sampling, restricting the
support of the generative distribution to trajectories admitting
coherent continuation across all active Postnikov stages.
No external ground-truth oracle is required at inference time.
Topology determines support; probability determines selection within
support.}
\end{quote}

\noindent
We are careful to maintain three precise distinctions:
\begin{itemize}
  \item Geometric obstruction $\neq$ semantic falsity.
        A blocked trajectory is unsupported by the compiled topology,
        not necessarily false in all possible worlds.
  \item Binary gate $\neq$ deterministic system.
        Within the admissible support, $P_\theta$ remains probabilistic.
  \item Structural constraint $\neq$ omniscience.
        The admissibility system encodes the compiled topology,
        not universal human knowledge.
\end{itemize}
The claim is:

\begin{quote}
\emph{A compiled semantic admissibility architecture in which
generation is transformed from pointwise likelihood maximisation
into sampling only from liftable continuations --- where
liftability is tested prior to softmax normalisation, at multiple
semantic scales, through compiled runtime admissibility projection
structures, requiring no external factual oracle at inference time.}
\end{quote}

\noindent
This is a nontrivial conceptual shift:
standard transformers \emph{maximise local likelihood};
this architecture \emph{samples only from liftable continuations}.
The empirical correlate is demonstrated in the Phase~1 benchmark
(Table~2: 100\% validity at all temperatures vs.\ baseline collapse
to 54\% at $T=3.0$).
The deepest remaining question --- whether the topology generalises
compositionally beyond the compiled training manifold --- is
the central programme of Phases~2--4.

The strongest formal statement is the Homological Admissibility
Theorem~\ref{thm:main}: any generated sequence whose continuation
induces a trivial obstruction class at any Postnikov stage is
excluded from decoding support --- algebraically, not
statistically.

\subsection*{Future work: compiler passes}

The immediate algebraic programme comprises:
\begin{itemize}
  \item \textbf{Pass 6} (Cone compiler): wall-crossing cluster
        mutations at interchange stations via SyzygyProjection.
  \item \textbf{Pass 7} (Vanishing cycles): seasonal Dehn twist
        monodromy via ExceptionalProjection.
  \item \textbf{Pass 8} (Stalk transport): Gauss-Manin connection
        for full $m_1$ parallel transport via SpectralProjection.
\end{itemize}

\subsection*{Future work: empirical validation requirements}

The framework is beyond idea stage.
To achieve publication, the following empirical evidence is required:

\begin{enumerate}[label=(\roman*)]
  \item \textbf{Benchmark comparison}: baseline transformer vs.\
        homologically-gated transformer on standard hallucination
        benchmarks (TruthfulQA, HaluEval, FactScore).
  \item \textbf{Invalid trajectory rate}: measure the fraction of
        generated sequences that are topologically blocked at each
        Postnikov stage. This quantifies how much work the gate
        is actually doing.
  \item \textbf{Long-context consistency}: demonstrate that the
        Floer-like continuation gate maintains coherence over
        longer contexts than the baseline.
  \item \textbf{OOD robustness}: show that the topology
        contributes nontrivially on out-of-distribution inputs
        --- not merely on the training distribution.
  \item \textbf{Ablation studies}: remove each Postnikov stage
        independently and measure the hallucination rate at that
        level. This validates the hierarchical structure.
  \item \textbf{Language model instantiation}: define the
        Lagrangians, symplectic form, and toric ideal for a
        concrete NLP task and validate the gate on factual QA.
\end{enumerate}

\noindent
Without these, reviewers will correctly characterise the system as
``a symbolic validator attached to a transformer.''
The experiments must prove that the topology contributes
nontrivially --- that the admissibility structure produces
compositional generalisation and emergent long-range consistency
that a standard learned mask cannot replicate.

% =============================================================================
\begin{thebibliography}{9}

\bibitem{FOOO}
  K.~Fukaya, Y.-G.~Oh, H.~Ohta, K.~Ono,
  \emph{Lagrangian Floer Theory and Mirror Symmetry on Compact
  Toric Manifolds},
  arXiv:1009.1648v3 [math.SG], 292 pp., 2016.

\bibitem{FLTZ}
  B.~Fang, C.-C.~M.~Liu, D.~Treumann, E.~Zaslow,
  \emph{T-duality and Homological Mirror Symmetry for Toric Varieties},
  Adv.\ Math.\ \textbf{229} (2012), 1873--1911.

\bibitem{Seidel}
  P.~Seidel,
  \emph{Fukaya Categories and Picard-Lefschetz Theory},
  European Mathematical Society, 2008.

\bibitem{4ti2}
  4ti2 team,
  \emph{4ti2 --- A software package for algebraic, geometric and
  combinatorial problems on linear spaces},
  \texttt{https://4ti2.github.io}.

\bibitem{Romer}
  P.~Romer,
  \emph{Mathiness in the Theory of Economic Growth},
  American Economic Review \textbf{105}(5) (2015), 89--93.

\end{thebibliography}

\end{document}
